% GENERATED FILE. Edit canonical lessons, then run npm run generate:books.
% canonical-source-hash: fnv1a64:3a66741fcab8ce8b
% canonical-lessons: HI-C354-vot, HI-C354-nagrik, HI-C354-janmat, HI-C354-loktantra, HI-C354-thopna

\chapter{Citizens and Votes}
\label{ch:hi-citizens-and-votes}
\hlchaptermodality{354}

\begin{chapteropening}
\textbf{By the end of this chapter:} \emph{I can talk about voting and the rights of citizens.}

Complete the last lesson of chapter 354: I can talk about voting and the rights of citizens.
\end{chapteropening}

\section[\texorpdfstring{\hi{वोट} (voṭ)}{वोट (voṭ)}]{\hi{वोट} (voṭ) — a vote}

\label{lesson:HI-C354-vot}



\begin{quote}
Before the new one: say the Hindi for a proposal, then the Hindi for to hold to be.
\end{quote}

\subsection*{You'll want to know: \hi{वोट}}
\textbf{\hi{वोट}} — \emph{voṭ} — \textquotedblleft{}a vote\textquotedblright{}.

\begin{quote}
\hi{मैंने} \hi{पिछले} \hi{चुनाव} \hi{में} \hi{वोट} \hi{दिया।}
\end{quote}

\begin{quote}
\emph{maiṁne pichle cunāv meṁ voṭ diyā.}
\end{quote}

\begin{quote}
I voted in the last election.
\end{quote}

Masculine.

\begin{cousinweb}
\textbf{\hi{वोट}} (\emph{voṭ}) is a loan from English \emph{vote}.
\end{cousinweb}

\begin{grammarlens}[title={What you've built}]
One more word for giving reasons and opinions.
\end{grammarlens}

\subsection*{Your turn}
\begin{itemize}
\raggedright
  \item \emph{Say it:} \emph{voṭ}
  \item \emph{Say it:} \emph{voṭ}, once more
  \item \emph{Say it:} \emph{maiṁne pichle cunāv meṁ voṭ diyā.} — I voted in the last election.
  \item \emph{Recall:} say the Hindi for a proposal, then the Hindi for to hold to be, then say \emph{voṭ} again
\end{itemize}

\begin{tcolorbox}[breakable,colback=teal!4,colframe=teal!35!black,title={Before you move on}]
What does \hi{वोट} mean? (\textquotedblleft{}A vote\textquotedblright{}.) Say it once more.
\end{tcolorbox}
\section[\texorpdfstring{\hi{नागरिक} (nāgrik)}{नागरिक (nāgrik)}]{\hi{नागरिक} (nāgrik) — a citizen}

\label{lesson:HI-C354-nagrik}



\begin{quote}
Before the new one: say the Hindi for to hold to be, then the Hindi for a vote.
\end{quote}

\subsection*{You'll want to know: \hi{नागरिक}}
\textbf{\hi{नागरिक}} — \emph{nāgrik} — \textquotedblleft{}a citizen\textquotedblright{}.

\begin{quote}
\hi{हर} \hi{नागरिक} \hi{को} \hi{वोट} \hi{देने} \hi{का} \hi{अधिकार} \hi{है।}
\end{quote}

\begin{quote}
\emph{har nāgrik ko voṭ dene kā adhikār hai.}
\end{quote}

\begin{quote}
Every citizen has the right to vote.
\end{quote}

Masculine or feminine, by the person meant.

\begin{cousinweb}
\textbf{\hi{नागरिक}} (\emph{nāgrik}) is a tatsama, borrowed from Sanskrit \emph{nāgarika}, \textquotedblleft{}a townsman; polite\textquotedblright{} (Monier-Williams).
\end{cousinweb}

\begin{grammarlens}[title={What you've built}]
One more word for giving reasons and opinions.
\end{grammarlens}

\subsection*{Your turn}
\begin{itemize}
\raggedright
  \item \emph{Say it:} \emph{nāgrik}
  \item \emph{Say it:} \emph{nāgrik}, once more
  \item \emph{Say it:} \emph{har nāgrik ko voṭ dene kā adhikār hai.} — Every citizen has the right to vote.
  \item \emph{Recall:} say the Hindi for to hold to be, then the Hindi for a vote, then say \emph{nāgrik} again
\end{itemize}

\begin{tcolorbox}[breakable,colback=teal!4,colframe=teal!35!black,title={Before you move on}]
What does \hi{नागरिक} mean? (\textquotedblleft{}A citizen\textquotedblright{}.) Say it once more.
\end{tcolorbox}
\section[\texorpdfstring{\hi{जनमत} (janmat)}{जनमत (janmat)}]{\hi{जनमत} (janmat) — public opinion}

\label{lesson:HI-C354-janmat}



\begin{quote}
Before the new one: say the Hindi for a vote, then the Hindi for a citizen.
\end{quote}

\subsection*{You'll want to know: \hi{जनमत}}
\textbf{\hi{जनमत}} — \emph{janmat} — \textquotedblleft{}public opinion\textquotedblright{}.

\begin{quote}
\hi{जनमत} \hi{नए} \hi{क़ानून} \hi{के} \hi{ख़िलाफ़} \hi{है।}
\end{quote}

\begin{quote}
\emph{janmat nae qānūn ke khilāf hai.}
\end{quote}

\begin{quote}
Public opinion is against the new law.
\end{quote}

Masculine.

\begin{cousinweb}
\textbf{\hi{जनमत}} (\emph{janmat}) joins two Sanskrit words, \emph{jana}, \textquotedblleft{}people\textquotedblright{}, and \emph{mata}, \textquotedblleft{}an opinion\textquotedblright{} (Monier-Williams): what the people think. The compound is a modern one.
\end{cousinweb}

\begin{grammarlens}[title={What you've built}]
One more word for giving reasons and opinions.
\end{grammarlens}

\subsection*{Your turn}
\begin{itemize}
\raggedright
  \item \emph{Say it:} \emph{janmat}
  \item \emph{Say it:} \emph{janmat}, once more
  \item \emph{Say it:} \emph{janmat nae qānūn ke khilāf hai.} — Public opinion is against the new law.
  \item \emph{Recall:} say the Hindi for a vote, then the Hindi for a citizen, then say \emph{janmat} again
\end{itemize}

\begin{tcolorbox}[breakable,colback=teal!4,colframe=teal!35!black,title={Before you move on}]
What does \hi{जनमत} mean? (\textquotedblleft{}Public opinion\textquotedblright{}.) Say it once more.
\end{tcolorbox}
\section[\texorpdfstring{\hi{लोकतंत्र} (loktantra)}{लोकतंत्र (loktantra)}]{\hi{लोकतंत्र} (loktantra) — democracy}

\label{lesson:HI-C354-loktantra}



\begin{quote}
Before the new one: say the Hindi for a citizen, then the Hindi for public opinion.
\end{quote}

\subsection*{You'll want to know: \hi{लोकतंत्र}}
\textbf{\hi{लोकतंत्र}} — \emph{loktantra} — \textquotedblleft{}democracy\textquotedblright{}.

\begin{quote}
\hi{लोकतंत्र} \hi{में} \hi{लोग} \hi{अपने} \hi{नेता} \hi{चुनते} \hi{हैं।}
\end{quote}

\begin{quote}
\emph{loktantra meṁ log apne netā cunte haiṁ.}
\end{quote}

\begin{quote}
In a democracy the people choose their leaders.
\end{quote}

Masculine.

\begin{cousinweb}
\textbf{\hi{लोकतंत्र}} (\emph{loktantra}) is a tatsama, borrowed from Sanskrit \emph{lokatantra}, \textquotedblleft{}the course of the world\textquotedblright{} (Monier-Williams). Modern Hindi gives the compound of \emph{loka}, \textquotedblleft{}people\textquotedblright{}, and \emph{tantra}, \textquotedblleft{}system\textquotedblright{}, a new sense: rule by the people.
\end{cousinweb}

\begin{grammarlens}[title={What you've built}]
One more word for giving reasons and opinions.
\end{grammarlens}

\subsection*{Your turn}
\begin{itemize}
\raggedright
  \item \emph{Say it:} \emph{loktantra}
  \item \emph{Say it:} \emph{loktantra}, once more
  \item \emph{Say it:} \emph{loktantra meṁ log apne netā cunte haiṁ.} — In a democracy the people choose their leaders.
  \item \emph{Recall:} say the Hindi for a citizen, then the Hindi for public opinion, then say \emph{loktantra} again
\end{itemize}

\begin{tcolorbox}[breakable,colback=teal!4,colframe=teal!35!black,title={Before you move on}]
What does \hi{लोकतंत्र} mean? (\textquotedblleft{}Democracy\textquotedblright{}.) Say it once more.
\end{tcolorbox}
\section[\texorpdfstring{\hi{थोपना} (thopnā)}{थोपना (thopnā)}]{\hi{थोपना} (thopnā) — to impose, to force (on someone)}

\label{lesson:HI-C354-thopna}



\begin{quote}
Before the new one: say the Hindi for public opinion, then the Hindi for democracy.
\end{quote}

\subsection*{You'll want to know: \hi{थोपना}}
\textbf{\hi{थोपना}} — \emph{thopnā} — \textquotedblleft{}to impose, to force (on someone)\textquotedblright{}.

\begin{quote}
\hi{कोई} \hi{तुम} \hi{पर} \hi{अपनी} \hi{राय} \hi{नहीं} \hi{थोप} \hi{सकता।}
\end{quote}

\begin{quote}
\emph{koī tum par apnī rāy nahīṃ thop saktā.}
\end{quote}

\begin{quote}
Nobody can force their opinion on you.
\end{quote}

Past, with \hi{ने} (he or she …): \hi{उसने} \hi{अपनी} \hi{राय} \hi{थोपी}, \emph{usne apnī rāy thopī}.

\begin{cousinweb}
\textbf{\hi{थोपना}} (\emph{thopnā}) is inherited from \emph{*stōpyate}, \textquotedblleft{}is obstructed\textquotedblright{}, a form CDIAL reconstructs from its descendants (CDIAL 13723). CDIAL glosses the Hindi verb \textquotedblleft{}to cover, to heap\textquotedblright{}; the sense of forcing a thing on someone is not among its glosses, and may have grown from heaping a thing on.
\end{cousinweb}

\begin{grammarlens}[title={What you've built}]
One more word for giving reasons and opinions.
\end{grammarlens}

\subsection*{Your turn}
\begin{itemize}
\raggedright
  \item \emph{Say it:} \emph{thopnā}
  \item \emph{Say it:} \emph{thopnā}, once more
  \item \emph{Say it:} \emph{koī tum par apnī rāy nahīṃ thop saktā.} — Nobody can force their opinion on you.
  \item \emph{Recall:} say the Hindi for public opinion, then the Hindi for democracy, then say \emph{thopnā} again
\end{itemize}

\begin{tcolorbox}[breakable,colback=teal!4,colframe=teal!35!black,title={Before you move on}]
What does \hi{थोपना} mean? (\textquotedblleft{}To impose, to force (on someone)\textquotedblright{}.) Say it once more.
\end{tcolorbox}
